\section{Discussion}\label{sec:discussion}

\subsection{The shape of the calculus}\label{sec:shape}

Table~\ref{tab:summary} collects the certificates, the response correspondence each one empties, and the design consequence each one licenses. The entries are of different kinds and the calculus depends on keeping them apart. The exit certificate is \emph{dynamic}: it defeats every control and needs no observation argument. The common-action obstruction is \emph{static-epistemic}: it defeats every policy of a held-action class at one information state, and Proposition~\ref{prop:emptiness} exhibits it at minimal size. The delayed-information obstruction interpolates between them in time, and the recourse obstruction detects the failure that only appears after information arrives. The fibre criterion limits \emph{verdicts}, not policies, and the certainty-equivalence trap limits one uncorrected policy. The finite recursion supplies the exact counterpart of all of them in finite systems.

\begin{table}[tb]
\centering\small\renewcommand{\arraystretch}{1.1}
\begin{tabularx}{\textwidth}{@{}>{\raggedright\arraybackslash}p{0.36\textwidth}>{\raggedright\arraybackslash}p{0.27\textwidth}X@{}}
\toprule
Certificate & Emptied response correspondence & Design consequence\\
\midrule
finite-time exit (Thm.~\ref{thm:exit}) & response set of some compatible state, under full information & none: dynamics, not information\\\addlinespace[3pt]
common action, admissibility and safety cases (Thm.~\ref{thm:common-action}; Prop.~\ref{prop:emptiness}; ladder, Prop.~\ref{prop:ladder}) & $U^B(B)$, $\cR^B_\cV(B)$, $\Atube(B,\Delta)$ & an observation that separates the incompatible states; a larger common command set\\\addlinespace[3pt]
delayed information (Thm.~\ref{thm:delayed}, Rem.~\ref{rem:sigma}, Thm.~\ref{thm:lp}) & blind-window survival: $\sigma^*_\Pi(B_0)<\Tobs$ & an earlier informative observation\\\addlinespace[3pt]
post-observation recourse (Prop.~\ref{prop:recourse}) & post-reveal continuation (oracle-relaxed) & more post-observation authority; a shorter blind window\\\addlinespace[3pt]
fibre certification (Prop.~\ref{prop:fibre}, Cor.~\ref{cor:certainly-safe}) & exact verdict: $K\ne O^{-1}(O(K))$ & a refined index; per-floor measurement; silence outside $\cY_{\mathrm{safe}}$\\\addlinespace[3pt]
certainty equivalence (Prop.~\ref{prop:ce-trap}) & kernel of the uncorrected controller & correct the known bias\\\addlinespace[3pt]
finite recursion (Thms.~\ref{thm:finite-horizon}, \ref{thm:onestep}) & $B_0\notin\cW_N$ & exact characterisation; obstruction tree\\\addlinespace[3pt]
\bottomrule
\end{tabularx}
\caption{The certificates, the response correspondence each empties, and the design consequence each licenses. The exit, common-action, timing and recourse rows concern nonviability in their stated systems and classes; the fibre row concerns verdicts; the certainty-equivalence row one controller.}
\label{tab:summary}
\end{table}

A failure of observation can thus be \emph{diagnosed} by which certificate applies. If the exit certificate holds, no monitoring design helps. If the common-action obstruction holds, the remedy is an observation that distinguishes the incompatible states, and refinement is never harmful (Proposition~\ref{prop:monotone}), though the theorem claims no benefit beyond removing this test and other mechanisms may still certify nonviability on a finer partition. If the timing bound is violated with slack, the remedy is a shorter review interval or a faster indicator. If the recourse certificate fires, the remedy is authority after the observation rather than information before it. If the fibre criterion fails, the remedy is an observation that separates the sides of the floor rather than a better threshold on the aggregate. In the one-step exact class the design converse is immediate: an observation structure is \emph{adequate} on a target exactly when every fibre meeting the target has a nonempty common safe action---a discrete common-action test, distinct from the static fibre criterion, for which pairwise intersection is necessary and, off convexity, not sufficient (Remark~\ref{rem:count}).

\subsection{Position relative to barrier certificates, estimation tubes and reachability}\label{sec:position}

Barrier certificates certify safety from a scalar function whose zero level set separates the unsafe region from all trajectories \citep{prajna2004,prajna2007}, with a converse under convex-duality conditions \citep{prajna2005} and necessary-and-sufficient hybrid characterisations \citep{maghenem2019}. Theorem~\ref{thm:exit} is the unsafety complement of that converse: a Dini-drift certificate of unsafety that needs no observation structure and is constructive in the enforcing disturbance and exit time. The remaining certificates act at the level of the information set, which barrier constructions---presupposing a state or a specified estimate for feedback---do not address: a search restricted to observation-only barriers is inconclusive about the existence of a common safe action, and a state-dependent barrier need not be constant on fibres merely because the controller is output-based.

The estimation-tube programme carries imperfect measurement to set-valued dynamics on an estimation space \citep{kurzhanski1997,cardaliaguet2007}. Where the reduction applies with matching policy and admissibility hypotheses it is exact for values, and its belief-state controls are common controls. The present theorems delimit its reach rather than contradict it: the obstruction calculus is the local certificate layer of that programme---the common-action obstruction certifies that the belief-state predecessor is empty, the timing bound quantifies when it fails for reasons of information arrival, the fibre criterion says when no observation-based verdict can be exact. Belief-space viability answers \emph{whether}; the obstruction calculus answers \emph{why not, and what to change}. A single worked system carrying the three methods side by side is in S3.5, and the expanded review of the sufficiency results against which the calculus is defined (Veliov's condition, the estimation-tube reduction, observer-and-buffer transfer, and the linear substitution alternative) is S5.

Hamilton--Jacobi reachability \citep{mitchell2005,margellos2011,bansal2017} indexes its value by a physical state and is, like viability theory proper, a sufficiency machine; even under perfect observation its grid solvers scale exponentially in the state dimension, and the learning-based solvers that lift that bound \citep{bansal2021} purchase scalability with approximation. The explicitly finite and polyhedral certificates here address a different verification task---record-constrained set-membership viability---and nothing is claimed about partial-observation reachability methods with their own information-state formulations \citep{chatterjee2009}.

\paragraph{Probabilistic counterpart.} The information-state semantics of this paper is set-membership. Under a stochastic transition and observation likelihood on finite spaces, the maximal survival probability $V_k(b)$ obeys the Smallwood--Sondik belief recursion with an absorbing unsafe state, and the common-action obstruction has a chance-constrained form: if every action loses at least mass $\delta$ from the belief in one step, then $V_k(b)\le1-\delta$ for every $k\ge1$ and the chance constraint at level $1-\varepsilon$ is infeasible for $\varepsilon<\delta$. In the deterministic-kernel limit (single-successor transitions and deterministic observations) $V_k(b)=1$ exactly when $\supp(b)\in\cW_k$, with a deficit of at least the smallest compatible mass otherwise; this bound does not follow from support equality for a genuinely stochastic kernel, which is governed by its transition probabilities. Statements, proofs, the worked instances and the hold-class value table of the audit grid are in S4.

\subsection{Consequences for monitoring and indicator design}\label{sec:design}

Five consequences follow for the design of monitoring in sustainability governance; each is class-relative where its certificate is.

\emph{Timing.} Theorem~\ref{thm:delayed} gives a conditional sufficient bound, not a general minimum frequency: to avoid this failure test the informative observation cannot arrive strictly after $\inf q/\varepsilon$, the worst-case time from the most-constrained compatible state to violation, and an earlier reading is not by itself a viability guarantee. The bound is computable from the drift estimate that a robustness analysis computes in any case.

\emph{Coarseness.} Theorem~\ref{thm:common-action} says that an indicator is exposed when some observation fibre crosses the safe-control partition of the state space. The observation must separate safety classes, not states; keeping every fibre within one class removes this exposure.

\emph{Aggregation.} Proposition~\ref{prop:fibre} applies verbatim to composite indices. An index whose fibres cross a floor cannot support exact floor certification; the certainly-safe set is the domain in which an index-based verdict is sound. The requirement is sharpest where the floor constrains the productive base rather than its output, since an aggregate of services can read at full level while the base regenerating them is reduced beneath it. This is a formal complement to the thesis, argued informally since \citet{martinezalier1998}, that weakly comparable values cannot be fused into one metric without loss; the quantifier-order separation between scalarised and coordinate-wise feasibility that underlies it is developed in \citet{abaee2026aggregate}.

\emph{Bias.} Proposition~\ref{prop:ce-trap} shows that monitoring has two obligations---bounding the observation error and supplying the correction---and that the second lies entirely at the governance structure's discretion: the indicator is published, the correction is not applied, and the safety set of that one uncorrected law empties although a corrected policy remains viable. The remedy is structural correction of the observation map, not a better sensor.

\emph{Institutions.} The hierarchy \eqref{eq:hierarchy} records that epistemic contraction is one of three causes of kernel shrinkage, the others being disturbances and institutional restriction. The institutional kernel $\IRViab_\cJ(\cV)$ combines them: an institution restricted in what it may observe \emph{and} in what it may command inherits both contractions. The certificates remain valid, and typically tighten, when the command set is further restricted, since restriction of the action correspondence can only empty kernels (Proposition~\ref{prop:monotone}). These are design conclusions from the certificates, not empirical findings; their empirical testing belongs to applied studies.

A caution on sensor models: no universal numerical sensor is posited. Thresholded and tolerance-band readings can each fail closed-posterior or approximation obligations under their exact conventions, and any regular-observation lemma must verify joint approximation of reachable states, admissible actions, readings and exact posteriors rather than infer posterior regularity from continuity of the raw sensor.

\subsection{Limitations}\label{sec:limitations}

(i) The continuous-time certificates are sufficient conditions; completeness is claimed only where backward recursion is exact (Section~\ref{sec:finite}) and in the static class (Theorem~\ref{thm:static}). The middle ground---neither a Veliov-type condition nor an obstruction certificate---is open, and continuous completeness would additionally require a measurable-selection theory beyond the explicit selections the certificates exhibit \citep{bertsekas1978}. (ii) The timing obstruction is blind to the recourse mode; the recourse certificate covers it soundly in the affine class under an oracle relaxation, and a complete continuous-time characterisation without that relaxation is open (Open Problem~\ref{op:dynamic}). (iii) The analytic certificates rest on explicit adverse-path hypotheses (\hyp{C3}, \hyp{D2}, \hyp{E2}); when such a premise is unverified the certificate does not certify the model in question, and neither closed graph nor lower semicontinuity of $D$ supplies it. (iv) The semantics is set-membership; the probabilistic counterpart is separate (S4). (v) The examples are finite-dimensional; delay systems require the retarded extension of the comparison argument \citep{hale1993}, and observation error and hidden parameters require their own information models.
